
\ifCLASSOPTIONcompsoc
\IEEEraisesectionheading{\section{Introduction}\label{sec:introduction}}
\else
\section{Introduction}
\label{sec:introduction}
\fi





\subsection{Brief intro/overview}
\IEEEPARstart{T}{his} should include context of the topic and diss as well as importance. Brief.


\subsection{Research question and aims}

The research question guiding this project is:
\begin{quote}
\textit{``\textbf{How can modern web application vulnerabilities be authentically
modelled and pedagogically structured within Capture-the-Flag
environments to enhance learning outcomes in cybersecurity education?}''}
\end{quote}
[BRIEFLY DISCUSS]

\subsubsection{Deliverables}
The project deliverables are organised into three tiers of increasing
scope and complexity, summarised in Table~\ref{tab:deliverables}.
These tiers also relate to increasing difficulty, where \textit{Basic}
deliverables include fairly accessible single-step CTFs, while
\textit{Intermediate} and \textit{Advanced} tiers ramp up in both
exploit complexity and engineering rigour. Collectively, the
deliverables address a gap identified in the literature: existing
educational CTF platforms typically offer narrow vulnerability coverage
limited to classical flaws such as SQL injection and XSS
\cite{Sasiwanit}, lack per-user flag personalisation beyond simple
randomisation, and do not account for the ease with
which LLMs can solve uniform, well-documented challenges
\cite{ji2025}. This project targets all three shortcomings within a
single web-specific, graded coursework context.


\begin{table}[h]
\centering
\caption{Project Deliverables by Tier}
\label{tab:deliverables}
\begin{tabular}{|p{1.5cm}|p{6.2cm}|}
\hline
\textbf{Tier} & \textbf{Deliverables} \\ \hline
Basic &
Design and implement introductory CTF challenges targeting
foundational web vulnerabilities (e.g.\ insecure session handling,
weak authentication mechanisms). Each challenge is a standalone
web application with a single intended exploit path and per-user
flag personalisation. \\ \hline
Intermediate &
Design and implement medium-to-hard CTF challenges featuring
multi-stage exploit chains, where solving one stage reveals the
information required for the next. These challenges will also target more
complex vulnerability classes (e.g.\ SQL injection with filter
bypass, DOM XSS with bot-based privilege escalation, SSTI with
WAF evasion, SSRF with internal service pivoting) and require
players to synthesise knowledge across multiple techniques within
a single challenge. \\ \hline
Advanced &
End-to-end and unit testing for each challenge to verify that
intended exploit paths function correctly and that unintended
solutions are mitigated.
Full Dockerisation of each challenge as a self-contained
Docker Compose stack for portable distribution.
Per-user flag generation and credential management via the
shared challenge-generation system.
LLM-based difficulty validation using multiple models to assess
challenge resistance to automated solving.
Human participant evaluation with feedback collection. \\ \hline
\end{tabular}
\end{table}

\subsubsection{Measures of success}

\paragraph{\textbf{Basic}} Each challenge functions correctly when deployed, the
intended exploit path leads to a valid per-user flag, and the
application behaves realistically (i.e.\ non-exploitable features
function as expected alongside the deliberate vulnerability).
\paragraph{\textbf{Intermediate}} Multi-stage challenges are solvable through
their intended exploit chains without requiring knowledge external
to the challenge environment. Each stage teaches a discrete skill,
and exhibits a bread-crumbing mechanism (where one exploit reveals information
for the next or increases in difficulty). Post-design vulnerability auditing
confirms that no trivial unintended solution bypasses the intended
chain.
\paragraph{\textbf{Advanced}} Automated tests pass for each challenge, confirming that exploit paths remain functional after code changes. Docker
Compose stacks build and run successfully on a clean machine without
manual configuration beyond providing a student identifier. Human
testers of varying skill levels can engage with the challenges and
provide structured feedback. LLM evaluation produces documented
transcripts showing each model's approach, with results analysed
for patterns in where automated solvers succeed or fail.

\subsection{Achievements - What did I achieve?}

